\section{한계}
\label{sec:limitations}

\paragraph{합성 사용자.} 모든 정량적 결과는 추정기가 가정하는 것과 동일한 Bradley--Terry 모델을 따라 답하는 합성 사용자를 사용하며,
이는 모델 기반 추정기에 유리하다. Table~\ref{tab:robust}는 답변 모델과 선형성을 완화하지만, 그 사용자들 역시
시뮬레이션된 것이다. 실제 사람의 보행 스타일 선호가 이 특징들에 대해 선형인지,
안정적인지, 사람마다 다른지는 사용자 연구가 보여주어야 한다. Sec.~\ref{sec:pilot}의 파일럿은 진행 중이며 여기서는
결과를 보고하지 않는다. 우리의 검정력 시뮬레이션에 따르면, 사람들 간의 차이가 중간 정도일 때 개인 선호와 공유 선호를 구분하려면
약 20명의 참가자가 필요하며, 사람들이 거의 비슷하다면 현실적으로 가능한 어떤 인원으로도 충분하지 않다.

\paragraph{선형 특징.} 다섯 개의 특징은 수작업으로 설계되었고 한 명의 충돌 대상 사람에 대해 평가된다. 집단 관련 선호
(줄이나 대화 중인 사람들 사이를 가로지르기)에는 추가 특징이 필요하다.

\paragraph{폐루프.} 보행자는 리플레이되며 로봇에 반응하지 않는다. 인지는 완벽하고, 로봇은 동역학이 없는 점으로 취급된다.
GC에서 로봇이 유발한 충돌은 오라클을 포함한 모든 정책에서 기록된 보행자의 경우보다 높게 유지되는데, 이는 안전
필터가 등속 예측에 의존하기 때문이다. 학습된 예측기 또는 불확실성을 고려하는 예측기와 반응하는 보행자는 향후 과제이다.

\paragraph{생성기 비교.} flow-matching 후보는 900개의 GC 장면에서만 이용 가능했으므로(3\,s 시계, 플래너의
윈도우 안에 있는 사람들), 비교에는 그 시계, 그 사람들, 그리고 네 개의 특징을 사용한다.
